% Part: proof-theory
% Chapter: normalization
% Section: normalization-thm

\documentclass[../../../include/open-logic-section]{subfiles}

\begin{document}

\olfileid{pt}{nor}{thm}

\olsection{సాధారణీకరణ ప్రమేయం}

\tetoken{ఒక నిరూపణ}{!!^a{proof}}లో కట్ ఖండాలు లేకపోతే అది \emph{సాధారణ} రూపంలో ఉందంటాం. అటువంటి నిరూపణకు స్థానమార్పు పరివర్తనగానీ తగ్గింపు పరివర్తనగానీ వర్తింపజేయలేము; అలాగే వాటిలో ఏదీ వర్తించకపోతే అది సాధారణ రూపంలోనే ఉంటుంది. \tetoken{ఒక నిరూపణను}{!!a{proof}} \emph{సాధారణీకరించడం} అంటే కట్‌లు మిగలకుండా పోయే వరకు స్థానమార్పు, తగ్గింపు పరివర్తనలను వర్తింపజేసి దాన్ని సాధారణ నిరూపణగా మార్చడం. ఇది ఎప్పుడూ సాధ్యమని చెప్పేదే \emph{సాధారణీకరణ ప్రమేయం}.

\begin{thm}\ollabel{thm:normalization} ఏ \Log{N1i}-\tetoken{నిరూపణ}{!!{proof}}~$\delta$ అయినా సాధారణీకరించబడుతుంది; అంటే కట్‌లు లేని \tetoken{ఒక నిరూపణగా}{!!a{proof}}~$\delta'$ మార్చవచ్చు.
\end{thm}

\begin{proof}
  $\delta$ యొక్క కట్ స్థాయి, కట్ పొడవుపై ఆగమనంతో నిరూపిద్దాం. ఆగమన ఉపపత్తి: కట్ స్థాయి తక్కువగా ఉన్న, లేదా అదే కట్ స్థాయి ఉండి కట్ పొడవు తక్కువగా ఉన్న ప్రతి \tetoken{నిరూపణ}{!!{proof}} సాధారణీకరించబడుతుంది.

  కట్ పొడవు~$0$ అయితే కట్‌లేవీ లేవు;~$\delta$ ఇప్పటికే సాధారణ రూపంలో ఉంది. కాబట్టి~$\delta$ యొక్క కట్ స్థాయి, కట్ పొడవు రెండూ~$>
  0$ అనుకుందాం. అత్యుపరి, అత్యంత కుడివైపు ఉన్న గరిష్ఠ కట్ ఖండాన్ని ఎంచుకుందాం. ఆ కట్ ఖండం పొడవు $>1$ అయితే స్థానమార్పు పరివర్తనను వర్తింపజేయండి. పొడవు~$1$ అయితే తగిన తగ్గింపు పరివర్తనను వర్తింపజేయండి. ఫలితంగా వచ్చే కొత్త \tetoken{నిరూపణ}{!!{proof}}కు అదే కట్ స్థాయి ఉండి కట్ పొడవు తక్కువగా ఉంటుంది, లేదా కట్ స్థాయి తక్కువగా ఉంటుంది; అప్పుడు ఆగమన ఉపపత్తి వర్తిస్తుంది.
\end{proof}

పై నిరూపణలో ఉపయోగించిన పద్ధతి—ఎల్లప్పుడూ అత్యుపరి, అత్యంత కుడివైపు ఉన్న గరిష్ఠ కట్‌ను పరిష్కరించడం—\tetoken{ఒక నిరూపణను}{!!a{proof}} సాధారణ రూపంలోకి తేవడానికి పరివర్తనలను వర్తింపజేయవలసిన ఒక నిర్దిష్ట క్రమాన్ని సూచిస్తుంది. ఆశ్చర్యకరంగా, వాస్తవానికి ఆ క్రమం ముఖ్యం కాదు. స్థానమార్పు, తగ్గింపు పరివర్తనలను ఏ క్రమంలో వర్తింపజేసినా చివరకు సాధారణ \tetoken{నిరూపణ}{!!{proof}} వస్తుంది.

ఖండం, కట్ నిర్వచనాలను \Log{N2i}కు సులభంగా రూపొందించవచ్చు; \Log{N1i}లోని స్థానమార్పు, తగ్గింపు పరివర్తనలు నేరుగా \Log{N2i}కు మారుతాయి. అందువల్ల సాధారణీకరణ ప్రమేయం \Log{N2i}కూ వర్తిస్తుంది.

\begin{thm}\ollabel{thm:normalization-N2i} ఏ \Log{N2i}-\tetoken{నిరూపణ}{!!{proof}}~$\delta$ అయినా సాధారణీకరించబడుతుంది; అంటే కట్‌లు లేని \tetoken{ఒక నిరూపణగా}{!!a{proof}}~$\delta'$ మార్చవచ్చు.
\end{thm}

\end{document}
